\documentclass[12pt]{article}
\usepackage{amsmath, amssymb}
\usepackage{geometry}
\usepackage{graphicx}
\geometry{margin=1in}

\title{\textbf{Lecture Scribe}\\Fundamentals of Probability in Computing}
\author{
Name: Sakshi Patel\\
Student ID: AU2440206\\
Course: Fundamentals of Probability in Computing\\
}

\begin{document}

\maketitle

\section*{�� PART 1: MAIN TOPICS (Roadmap)}

From your lecture, the core topics are:
\begin{itemize}
\item Why Joint Random Variables
\item Discrete Case – Joint PMF
\item Joint PMF Table Representation
\item Example: Rolling Two Dice
\item Continuous Case – Joint PDF
\item Geometric Interpretation
\item Joint Cumulative Distribution Function (CDF)
\item Continuous Case
\item Discrete Case
\item Properties of Joint CDF
\item Comprehensive Examples
\end{itemize}

\section*{�� 1. WHY JOINT RANDOM VARIABLES?}

\textbf{Definition}

A joint random variable studies two or more random variables together, instead of individually.

\textbf{Conceptual Understanding}

Earlier: You studied $X$ alone → like speed

Now: You study $(X, Y)$ → speed AND traffic density

\textbf{Why is this needed?}

Because real-world systems are connected, not isolated.

From your lecture:
\begin{itemize}
\item Speed depends on traffic
\item Throughput depends on SNR
\item Weather depends on multiple factors
\end{itemize}

\textbf{Key Insight}

Two variables can be:
\begin{itemize}
\item Independent → no relation
\item Dependent → one affects the other
\end{itemize}

\section*{�� 2. DISCRETE CASE — JOINT PMF}

\textbf{Definition}

\[
P(X = x, Y = y)
\]

It gives probability that both events happen together.

\textbf{Properties}

\[
P(x,y) \geq 0
\]

\[
\sum_x \sum_y P(x,y) = 1
\]

\section*{�� 3. JOINT PMF TABLE}

\textbf{Representation}

\begin{center}
\begin{tabular}{c|ccc}
 & $Y=1$ & $Y=2$ & $Y=3$ \\
\hline
$X=1$ & $p_{11}$ & $p_{12}$ & $p_{13}$ \\
$X=2$ & $p_{21}$ & $p_{22}$ & $p_{23}$ \\
\end{tabular}
\end{center}

\section*{�� 4. EXAMPLE: TWO DICE}

Let:
\[
X = \text{Dice 1}, \quad Y = \text{Dice 2}
\]

Total outcomes:
\[
6 \times 6 = 36
\]

\[
P(X=x, Y=y) = \frac{1}{36}
\]

\textbf{Example Problem}

Find:
\[
P(X + Y = 7)
\]

Valid pairs:
\[
(1,6), (2,5), (3,4), (4,3), (5,2), (6,1)
\]

\[
P = \frac{6}{36} = \frac{1}{6}
\]

\section*{�� 5. CONTINUOUS CASE — JOINT PDF}

\textbf{Definition}

\[
f(x,y)
\]

\[
P(X=x, Y=y) = 0
\]

\textbf{Probability Formula}

\[
P(a \leq X \leq b, c \leq Y \leq d)
= \int_c^d \int_a^b f(x,y)\,dx\,dy
\]

\section*{�� 6. GEOMETRIC INTERPRETATION}

X-axis → $X$ \\
Y-axis → $Y$ \\
Height → $f(x,y)$

Probability = volume under surface

\section*{�� 7. CONTINUOUS EXAMPLE}

\[
f(x,y) = k \quad \text{for } 0 \leq x \leq 1, 0 \leq y \leq 1
\]

\[
\int_0^1 \int_0^1 k\,dx\,dy = 1
\]

\[
k(1 \times 1) = 1 \Rightarrow k = 1
\]

\section*{�� 8. JOINT CDF}

\textbf{Definition}

\[
F(x,y) = P(X \leq x, Y \leq y)
\]

\textbf{Continuous Case}

\[
F(x,y) = \int_{-\infty}^{x} \int_{-\infty}^{y} f(u,v)\,du\,dv
\]

\textbf{Discrete Case}

\[
F(x,y) = \sum_{u \leq x} \sum_{v \leq y} P(u,v)
\]

\section*{⚙️ 9. PROPERTIES OF JOINT CDF}

\[
0 \leq F(x,y) \leq 1
\]

Non-decreasing

Limits:
\begin{itemize}
\item $\to 1$ as $x,y \to \infty$
\item $\to 0$ as $x,y \to -\infty$
\end{itemize}

\section*{�� 10. CONNECTION BETWEEN PMF, PDF, CDF}

\begin{center}
\begin{tabular}{c|c}
Concept & Meaning \\
\hline
PMF & Exact probability (discrete) \\
PDF & Density (continuous) \\
CDF & Accumulated probability \\
\end{tabular}
\end{center}

\section*{�� FINAL SUMMARY}

\begin{itemize}
\item Joint RVs study multiple variables together
\item PMF → discrete probability
\item PDF → continuous density
\item CDF → cumulative probability
\item Discrete → counting
\item Continuous → integration
\end{itemize}

\section*{�� KEY FORMULA SHEET}

\textbf{Joint PMF}
\[
P(X=x, Y=y)
\]

\textbf{Total Probability}
\[
\sum_x \sum_y P(x,y) = 1
\]

\textbf{Joint PDF}
\[
f(x,y)
\]

\textbf{Probability (continuous)}
\[
\int \int f(x,y)\,dx\,dy
\]

\textbf{Joint CDF}
\[
F(x,y) = P(X \leq x, Y \leq y)
\]

\section*{�� PROBABLE EXAM QUESTIONS}

\textbf{High Probability}
\begin{itemize}
\item Define joint PMF, PDF, CDF
\item Dice-based problems
\item Compute probabilities from table
\end{itemize}

\textbf{Medium}
\begin{itemize}
\item Find constant $k$ in PDF
\item Compute region probability
\end{itemize}

\textbf{Conceptual Traps}
\begin{itemize}
\item Difference between PMF and PDF
\item Why PDF at a point = 0
\item Valid PDF conditions
\end{itemize}

\section*{�� FINAL TUTOR ADVICE}

If you don’t visualize, you won’t survive probability.

Always ask:
\begin{itemize}
\item “Am I counting outcomes?” → discrete
\item “Am I measuring area?” → continuous
\end{itemize}

\end{document}
